\section{Evaluation Protocol}
\label{sec:setup}

\subsection{Cases and Methods}
\label{sec:setup-cases}
The 68 cases combine the two Kusiak--Song wind data sets~\cite{Kusiak2010,Bansal2017} (Table~\ref{S-tab:S-wind}) with circular farms of radii $r=500$, 750 and 1000~m and $N=2,\ldots,10$, $2,\ldots,12$ and $2,\ldots,15$ turbines ($R=38.5$~m, $\ell_{\min}=4D=308$~m, $C_T=0.8$, $K=0.075$).

The \emph{main comparison} includes eight methods ($N_p=30$ where applicable): PSO-VNS (Algorithm~\ref{alg:hybrid}, $\omega=0.5$); PSO; SSA-VNS; SSA; LX-SSA; DE/rand/1/bin~\cite{Storn1997} (scale factor 0.5, crossover rate 0.9); VNS; and MS-SLSQP, a multistart sequential least-squares programming method (SciPy's SLSQP~\cite{Kraft1988}, forward-difference gradients) that minimizes the wake loss under explicit constraints, restarting from the initial, then random layouts (Section~\ref{S-sec:S-protocol}); all others minimize $F_p$~\eqref{eq:penalty} over $[-r,r]^{2N}$. The \emph{component analysis} (Section~\ref{sec:ablation}) adds LX-SSA-VNS and two random-sampling controls (Section~\ref{sec:template}; RSD-VNS on the benchmark only): after the common initial population, \emph{RSD-VNS} samples Phase~1 uniformly in the farm disc and \emph{RS-VNS} in the bounding square, 3,015 layouts versus the hybrids' 3,030 Phase-1 evaluations. A split study varies $\omega$ on \NSplitCases{} cases (Section~\ref{sec:split}). Table~\ref{tab:design} lists all studies of the paper with their cases, methods, budgets, seeds and runs; the re-evaluations and diagnostics use stored layouts and runs.
% CHECK-FINAL [C61]: RS-VNS and RSD-VNS switch after round(0.5 x 6,030) = 3,015 Phase-1 evaluations; RSD-VNS = experiment rsdisc (68 cases x 30 seeds only; not in the budget, feasible-start, Horns Rev or IEA37 studies); CHECK-DIAG [D17]: RS-VNS Phase-1 replay

\begin{table}[!htbp]
\centering
\caption{Experimental design. Budget: objective evaluations per run, all counted (initial population, finite-difference evaluations of MS-SLSQP); r: random, f: feasibility-preserving initialization. Runs: new optimization runs (``reused'': taken from another study). Ten methods: the eight of the main comparison, LX-SSA-VNS and RS-VNS. Split cases: middle and largest $N$ per farm and data set; six largest: largest $N$ per farm and data set. In total, 36{,}190 optimization runs.}
\label{tab:design}
\scriptsize
\setlength{\tabcolsep}{3pt}
\begin{tabular}{@{}>{\raggedright\arraybackslash}p{2.0cm}>{\raggedright\arraybackslash}p{1.85cm}>{\raggedright\arraybackslash}p{5.6cm}>{\raggedright\arraybackslash}p{2.15cm}>{\raggedright\arraybackslash}p{1.05cm}>{\raggedright\arraybackslash}p{1.15cm}>{\raggedright\arraybackslash}p{0.55cm}@{}}
\toprule
Study & Cases & Methods & Budget (init.) & Seeds & Runs & Sec. \\
\midrule
Main comparison & 68 benchmark & PSO-VNS, PSO, SSA-VNS, SSA, LX-SSA, DE, VNS, MS-SLSQP & 6,030 (r) & 1--30 & \NRunsMain & \ref{sec:overall} \\
Component analysis & 68 & LX-SSA-VNS, RS-VNS, RSD-VNS (\NDRsDiscRuns{} each); six reused & 6,030 (r) & 1--30 & 6{,}120 & \ref{sec:ablation} \\
Old PSO setting & 68 & PSO, $w=0.7$, $c_1=c_2=2$ & 6,030 (r) & 1--30 & \NDStoredRunsOld & \ref{sec:psosetting} \\
Budget split & \NSCases{} split cases & PSO-VNS, $\omega=0.25$, 0.75, 0.9 (\NDOmegaNinetyRuns{} each); $\omega=0.5$, 1 reused & 6,030 (r) & 1--30 & 1{,}080 & \ref{sec:split} \\
Coefficient sweep & \NSCases{} split cases & PSO, $w=0.7$, $c_1=c_2=c=1.2$, 1.4, 1.6--2.0 (step 0.1); $c=2$ with $v\gets0$ or $\lvert v\rvert\le\NSVmaxFrac\,(u-l)$ (\NSRuns{} each) & 6,030 (r) & 1--30 & 3{,}240 & \ref{sec:psosetting} \\
Budget scaling & six largest & ten methods (6,030 reused) & 30,030; 120,030 (r) & 1--30 & 3{,}600 & \ref{sec:budget} \\
Feasible starts & six largest & ten methods except RS-VNS & 6,030 (f) & 1--30 & 1{,}620 & \ref{sec:feasinit} \\
Horns Rev~1 & 16 turbines & ten methods (f: except RS-VNS) & 6,030 (r, f); 30,030, 120,030 (r) & 1--30; 1--10 at 120,030 & \NFHRNRuns & \ref{sec:hornsrev-site} \\
IEA37 CS~1 & 16, 36 turbines & ten methods & 6,030; 30,030 (r) & 1--30 & 1{,}200 & \ref{sec:iea37} \\
\midrule
Re-evaluation & stored final layouts & cubic power curve, Gaussian wake (main comparison); 5\textdegree{} and 1\textdegree{} bins: \NFNLayouts{} layouts (11 methods), \NFHRNRuns{} Horns Rev~1 runs (also PyWake NOJ) & -- & -- & -- & \ref{sec:robustness} \\
Diagnostics & \NDDynCases; \NDVnsCases; \NDFsCases & instrumented: PSO (both settings); Phase~2 of PSO-VNS, SSA-VNS, RS-VNS; PSO, DE from feasible starts & as stored & \NDDynSeeds; \NDVnsSeeds; \NDFsSeeds & \NDVerRuns{} (reruns) & \ref{S-sec:S-diag} \\
\bottomrule
\end{tabular}
\end{table}
% CHECK (manual, 2026-09-29; run counts derived from the design constants and confirmed by the row counts of the run files in analysis/, no macro unless named):
%   main 8 x 68 x 30 = 16,320 = \NRunsMain; component analysis 3 x 68 x 30 = 6,120 new (LXBV in fresh_bgrid, mpce_rsvns, mpce_rsdisc = \NDRsDiscRuns), \NRunsBenchmark = 16,320 + 6,120 = 11 x 2,040;
%   split: mpce_psosplit 720 rows (omega 0.25 / 0.75) + mpce_omega90 360 rows (\NDOmegaNinetyRuns) = 1,080; csweep: mpce_csweep_s0of1.csv 3,240 rows = 9 settings x \NSRuns;
%   budget: mpce_b30k + b30kp 2,100 rows and mpce_b120k + b120kp 1,900 rows, minus the superseded Horns Rev rows (300 and 100; replaced by hrfix) = 1,800 + 1,800 = 3,600 (10 x 6 x 30 x 2);
%   feasible init.: mpce_feas 1,680 + mpce_feasp 210 = 1,890 rows minus 270 superseded Horns Rev rows = 1,620 (9 x 6 x 30);
%   Horns Rev: mpce_hrfix 970 rows = 10 x 30 (6,030) + 10 x 30 (30,030) + 10 x 10 (120,030) + 9 x 30 (feasible) = \NFHRNRuns;
%   IEA37: mpce_iea16/iea36(+p) 1,200 rows = 10 methods x 2 scenarios x 2 budgets x 30 (Table S-tab:iea37-detail lists all ten; the main-text Table tab:iea37 shows eight);
%   total 22,440 + 2,040 + 1,080 + 3,240 + 3,600 + 1,620 + 970 + 1,200 = 36,190;
%   diagnostics \NDVerRuns = 2 x \NDDynRuns + 3 x \NDVnsRunsPSOVNS + \NDFsRunsPSO + \NDFsRunsDE (120 + 360 + 42); CHECK-DIAG [D01], [D17]

\subsection{Comparison Design and Statistics}
\label{sec:setup-stats}
\emph{Equal budgets and paired seeds.} Unless stated otherwise, a run uses exactly $B=6{,}030$ objective evaluations, including the initial population and MS-SLSQP's gradient evaluations. Each method uses seeds 1--30 and draws its initial population first, so runs with equal seeds share it and are paired.
% verified 2026-09-28: first Curve value identical across fresh_grid/vgrid/bgrid and mpce_psoc/psobv/rsvns/slsqp (13 files) for all 792 case-seed pairs with a feasible initial best (else NaN)

\emph{Feasibility-aware ranking.} Let $f_{a,c}$ be the number of feasible runs of method $a$ in case $c$ (tolerance $10^{-6}$~m; Section~\ref{S-sec:S-tol}) and $\bar L_{a,c}$ the mean wake loss of these runs. Method $a$ is \emph{qualified} in case $c$ if at least half of its runs are feasible (15 of 30, 5 of 10 for 10 seeds). Within a case, the qualified methods are ranked by the mean benchmark objective of their feasible runs (rounded to $10^{-6}$, so that exact ties, e.g., at the wake-free optimum for $N=2$, share their average rank), and the others below them, by their number of feasible runs. This limits, but does not remove, survivorship bias.

\emph{Tests.} All tests are two-sided ($\alpha=0.05$), with Holm's correction~\cite{Holm1979} within each family of comparisons.
1)~\emph{Case ranks}: the Friedman test on the case-by-method rank matrix with the Iman--Davenport $F$ statistic, and post hoc $z$ tests of the average rank $\bar R_a$ of every method against that of PSO-VNS, $z=(\bar R_a-\bar R_{\text{PSO-VNS}})/\sqrt{k(k+1)/(6n)}$ for $k$ methods and $n$ cases~\cite{Demsar2006,Derrac2011,LaTorre2021}.
2)~\emph{Case means} (primary, as mean-rank post hoc tests depend on the pool~\cite{Benavoli2016}): for methods $A$ and $B$, the case-mean difference is $d_c=\bar L_{A,c}-\bar L_{B,c}$ (percentage points, pp; negative: $A$ better), and the Wilcoxon signed-rank test is applied to the $d_c$ of each compared pair (in the main comparison, PSO-VNS against each method); a case in which only one of the two methods is qualified counts as a maximal difference in its favor (larger than every observed difference; cases: \NCmImputedPhrase), and one in which neither is qualified is dropped (\NCmDropped{} cases). The mean $\overline{\Delta L}$ of the $d_c$ over the cases in which both are qualified carries a 95\% percentile bootstrap CI (10{,}000 resamples of the cases, fixed seed). As the estimate that belongs to the Wilcoxon test we also give, for key pairs, the Hodges--Lehmann estimate, the median of the Walsh averages $(d_c+d_{c'})/2$, $c\le c'$, with the 95\% CI from the exact inversion of the test (Table~\ref{S-tab:X-hl}); it can differ from $\overline{\Delta L}$ when a few cases with large differences dominate the mean.
3)~\emph{Runs}: per case, Wilcoxon signed-rank tests on the 30 seed-paired runs, scored feasibility first (infeasible runs below every feasible run and ordered by minimum spacing; differences $\lvert d\rvert\le10^{-9}$ dropped), with the rank-biserial correlation $r_{\rm rb}$ and the counts of cases with a significant win, no significant difference and a significant loss (W/T/L, descriptive tallies).
\emph{Holm families.} A family is the set of comparisons of one table: for the rank post hoc tests, the $k-1$ comparisons with PSO-VNS of one Friedman analysis; for the case means, the comparisons of PSO-VNS with the other seven methods of the main comparison, or the \NAblNContrasts{} contrasts of the component analysis (Table~\ref{tab:ablation}); for the run-level tests, the comparisons of one table within one case. The main-comparison and the component-analysis tables are different families, so the same pair can have slightly different W/T/L tallies in the two; a single family over all \NXMultN{} case-mean tests is examined in the robustness analysis below.
% CHECK (manual, 2026-09-29): constants in analysis/mpce_results.py: BOOT_N = 10000 (fixed seed), wil() drops |d| <= 1e-9, rank_rule rounds to 1e-6 and ranks non-qualified methods by nfeas (average ranks for ties); Holm families: friedman_block (vs focus), case_mean_wilcoxon (per table), paired_vs / ablation contrasts (per case and table)

\emph{Practical equivalence.} A nonsignificant test does not show that two methods are equally good, and a significant difference can be negligible; we therefore test equivalence with two one-sided tests (TOST)~\cite{Lakens2017}.

\begin{definition}[Practical equivalence at three levels of inference]
\label{def:S-equiv}
Let $A$ and $B$ be two methods, $d_c$ their case-mean difference over the $n$ cases in which both are qualified, $\overline{\Delta L}=\frac1n\sum_c d_c$, and $m>0$ a margin. $A$ and $B$ are \emph{practically equivalent at margin $m$} at a given level if both one-sided hypotheses $H_0^-\colon\Delta\le-m$ and $H_0^+\colon\Delta\ge m$ about the estimand $\Delta$ of that level are rejected at $\alpha=0.05$, i.e., if the 90\% interval for $\Delta$ lies inside $(-m,m)$. The three levels are:
\begin{enumerate}
\item[(a)] \emph{Seed level} (the fixed benchmark). Estimand: the benchmark average of the expected case-mean differences of the 68 given cases; only the run-to-run variability is random. Interval: percentile bootstrap that resamples the 30 seed-paired runs within every case (the same seeds for both methods; case means over the feasible resampled runs).
\item[(b)] \emph{Case level} (exchangeable cases). Estimand: the expected case-mean difference over a population of cases like the 68; the cases are treated as independent draws. Interval: percentile bootstrap over the cases.
\item[(c)] \emph{Cluster level} (new farms). Estimand: the expected difference over new (data set, radius) clusters; the 68 cases form \NXClusters{} clusters with nested $N$. Interval: cluster-robust $t$ interval with the bias-reduced CR2 variance and a $t$ distribution with $G-1=5$ degrees of freedom, and a restricted wild-cluster bootstrap-$t$ TOST (CR2-studentized, all $6^6$ draws of Webb's six-point weights enumerated); equivalence at this level requires both.
\end{enumerate}
The \emph{minimal margin} of a level, $m_{\min}=\max(\lvert\ell\rvert,\lvert u\rvert)$ for the 90\% interval $[\ell,u]$, is the smallest margin at which equivalence holds: it holds for every $m>m_{\min}$ and for none smaller.
\end{definition}
TOST has size at most $\alpha$ without a multiplicity correction (intersection--union principle; Section~\ref{S-sec:S-th-equiv}). The estimate $\overline{\Delta L}$ is the same at every level, but the interval widens from the seed to the cluster level, so $m_{\min}$ grows and an equivalence statement must name its level (Table~\ref{tab:X-equiv-levels}, Fig.~\ref{fig:equiv-curve}). ``$A$ gives no gain larger than $m$ over $B$'' (non-superiority) needs only one of the two tests: the lower end of the 90\% interval must lie above $-m$. A result that is neither equivalent nor significantly different is inconclusive and supports no statement of the kind ``as good as''. The margin $m=\NEqMargin$~pp (\NXMarginMWhTurbDsI{} and \NXMarginMWhTurbDsII~MWh/yr per turbine in Data Sets I and II) is a post hoc convention; the minimal margins $m_{\min}$ let readers apply their own.

\emph{Bayesian signed-rank test.} The Bayesian signed-rank test~\cite{Benavoli2017} places a Dirichlet-process prior (strength 0.5, pseudo-observation at 0) on the distribution of the $d_c$ and, for every posterior sample, computes the shares $\theta_A$, $\theta_{\rm rope}$ and $\theta_B$ of the Walsh averages $(d_c+d_{c'})/2$ below $-m$, inside the region of practical equivalence (ROPE) $[-m,m]$ and above $m$. We report the posterior probability that equivalence is the most probable of the three outcomes, the share of the 50{,}000 posterior samples in which $\theta_{\rm rope}$ is the largest (Tables~\ref{tab:equivalence} and~\ref{S-tab:X-bayes}). It is not the probability that $\Delta$ lies within $\pm m$; it concerns the typical pairwise-averaged case difference, so it can disagree with TOST, and it treats the cases as exchangeable.
% CHECK-FINAL [C49]-[C52]: margin_pp == EQ_MARGIN; TOST: 90 % bootstrap CI inside (-m, m); Bayes signed-rank, rope (-m, m); min_margin_pp in tab:equivalence
% CHECK-EXTRA [X38]-[X42]: seed / case / cluster (CR2 t(5), wild-cluster Webb) levels and minimal margins (tab:X-equiv-levels); [X56]: margin = \NXMarginMWhTurbDsI / \NXMarginMWhTurbDsII MWh/yr per turbine; [X48]: Bayes = P(rope most probable); [X49]-[X51]: Hodges-Lehmann (tab:X-hl). Model-shift margin justification (X29/X30) deliberately dropped (D13).
% CHECK (manual, 2026-09-29): analysis/mpce_inference_extra.py: seed level resamples the 30 seed indices within every case (paired), cr2() uses t with G - 1 = 5 d.f., webb_all() enumerates 6^6 = 46,656 Webb draws, cluster-level equivalence = CR2 CI inside +-m and wild p_TOST <= 0.05; mpce_results.py: tost() add-one corrected tail shares, BAYES_S = 0.5, BAYES_Z0 = 0, BAYES_N = 50000. Definition adapted from optA/supp_theory.tex (def:S-equiv moved to the main text, see optA/sw/moved.md)

\emph{Robustness of inference.} The 68 cases form \NXClusters{} (data set, radius) clusters with nested $N$. Six qualification thresholds (\NXThrList{} feasible runs) change no verdict; cluster analyses (sign per cluster, cluster-robust $t$ with CR1 variance and 5 degrees of freedom, cluster bootstrap, leave-one-out) qualify only the significant contrasts of the salp-swarm hybrids with the controls; one Holm family over all \NXMultN{} case-mean tests makes only \NXMultLostHolm{} nonsignificant (Tables~\ref{S-tab:X-threshold}--\ref{S-tab:X-multiplicity}); Section~\ref{sec:robustness} adds a 1\textdegree{} direction resolution. As the cases are not independent, all $p$ values of the case-level tests describe this benchmark, not a population of farms.
% CHECK-EXTRA [X02]-[X09]/[X31]: six thresholds (NXThrList) change no verdict; [X10]: six clusters; [X11]: significant pairs unanimous in 6/6 clusters except \NXClCaseSigNotUnanimous (the significant contrasts of the salp-swarm hybrids with the controls are exactly these two plus SSA-VNS vs RS-VNS; LX-SSA-VNS vs RS-VNS is not significant, [X13]); [X16]/[X17]: cluster-robust t not significant for the two RSD-VNS contrasts and SSA-VNS vs RS-VNS (p \NXClSSAVNSvsRSVNSCRP); [X14]: LOCO changes only \NXLocoChanged; [X18]: one Holm family over the \NXMultN case-mean Wilcoxon tests loses \NXMultLostHolm. CHECK-DIR [F01]: 1-deg sub-bin re-evaluation (one sub-bin reproduces the benchmark objective)
% CHECK (manual, 2026-09-29): tab:X-cluster caption (mpce_supp_inference.tex): CR = cluster-robust t of the case-weighted mean, CR1 variance, 5 d.f.

\subsection{Tests Beyond the Benchmark}
\label{sec:setup-beyond}
\emph{Initialization and budget.} The feasible initialization draws seeded layouts inside the site and separates the turbines with L-BFGS-B on a packing penalty, without objective evaluations (Section~\ref{S-sec:S-protocol}); all methods but RS-VNS use it at 6,030 evaluations. With random starts, all ten run at 6,030, 30,030 and 120,030 evaluations. Both studies use the six largest cases (largest $N$ per farm and data set) and the Horns Rev~1 block, with 30 seeds (10 for Horns Rev~1 at 120,030).

\emph{Other sites and models.} The Horns Rev~1 block is PyWake's 16-turbine example~\cite{PyWake}: V80 turbines ($D=80$~m), the measured 12-sector wind climate, the parallelogram through the corner turbines as boundary (Section~\ref{sec:hornsrev-site}), $\ell_{\min}=4D$ and the Jensen wake with $K=0.04$ and $C_T$ at the free-stream speed.
% CHECK (manual, resolved 2026-09-29): description matches analysis/hornsrev_model.py (V80, 12 sectors, K = 0.04, hub-centre top-hat, C_T at free-stream speed, 5-deg bins after the hrfix) and supplement sec:S-hrmodel; all HR runs are the hrfix reruns
The 16- and 36-turbine scenarios of IEA37 Case Study~1~\cite{Baker2019,IEA37repo} (circular sites, radii 1300 and 2000~m, $\ell_{\min}=2D$, 16 directions at one speed, a simplified Bastankhah Gaussian wake) test optimizers rather than represent a site; the ten methods of the budget study run on them at 6,030 and 30,030 evaluations (30 seeds; Table~\ref{tab:iea37} shows the eight of the main comparison).
% CHECK (manual, 2026-09-29): mpce_iea16/iea36(+p) contain PSOBV, PSOC, SSABV, SSA, LXSSA, DE, BVNS, SLSQP, LXBV, RSVNS (120 runs each); formerly "the eight methods of the main comparison run on them"

\subsection{Verification and Reproducibility}
\label{sec:validation}
\emph{Automated claim checks.} The results quoted in the text and the tables are macros and tables generated by the analysis scripts from the stored per-run results. Every statement that depends on an outcome (a sign, a significance or equivalence verdict, an ordering, a count) is tied in the source to a condition that a check script evaluates on the stored summaries; the build runs these scripts after regenerating the results. There are 62 conditions for the main results, 57 for the statistical robustness analyses, 22 for the diagnostics, 41 for the direction-resolution re-evaluation and 15 for the coefficient sweep, and 57 theory checks recompute every number stated in the propositions and proofs (Section~\ref{S-sec:S-theory}) and confirm that it appears in the text. All pass for the version reported here.
% CHECK (manual, 2026-09-29; typed counts, no macro): analysis/check_final.log "62 PASS, 0 FAIL, 0 PENDING, 0 TAG ERRORS"; check_extra.log "57 PASS / 0 FAIL"; check_diag.log "22 PASS / 0 FAIL / 0 PENDING"; check_dir.log "41/41 checks passed"; check_csweep.log "15 PASS / 0 FAIL / 0 PENDING"; make_theory_figures.py --check "57 PASS, 0 FAIL" (optA/reviews2/R5_claims_audit.md; re-run on a scratch copy 2026-09-29).

\emph{Instrumented diagnostics.} The diagnostics (Table~\ref{tab:design}) rerun stored runs with instrumented copies of the optimizers, which differ from the code of the paper only by logging lines that read the state and draw no random numbers. All instrumented runs reproduce the stored runs, with the same wake loss, coordinates, convergence curve and number of evaluations (\NDVerStoredMatch{} of \NDVerRuns), and seed~1 of every case and method, rerun with the original, uninstrumented code in the same process, is bit-identical (\NDVerBitMatch{} of \NDVerBitRuns{} runs). Likewise, the point $c=2.0$ of the coefficient sweep reproduces the stored runs of the old setting bit for bit, and a replay of the random stream of Phase~1 of RS-VNS agrees with the stored runs (\NDRsReplayAgree{} of \NDRsReplayChecked).
% CHECK-DIAG [D01]: instrumented runs reproduce the stored runs (\NDVerStoredMatch/\NDVerRuns); [D02]: seed-1 reruns with the original classes bit-identical (\NDVerBitMatch/\NDVerBitRuns); [D17]: RS-VNS Phase-1 replay agrees; CHECK-CSWEEP [S01]: c = 2.0 reproduces the 360 stored old-setting runs bit for bit

\emph{Evaluator and code validation.} Our objective matches the original implementation (largest difference $8.7\times10^{-11}$ over 720 archived final objectives) and the official IEA37 calculator (relative difference ${<}10^{-11}$); the Horns Rev~1 model was compared with PyWake's NOJ model at identical bins (Section~\ref{sec:hornsrev-site}). Our reused SSA, LX-SSA, PSO and DE code reproduces the recorded result distributions of the original code (2 of 24 Mann--Whitney tests reject at $\alpha=0.05$; 1.2 expected).
% CHECK (manual, re-run 2026-09-29; no macro): analysis/validate_evaluator.py -> 720 records, max|dObj| 8.73e-11; python3 analysis/iea37_model.py -> rel.err 8.8e-12 / 2.8e-12 / 2.9e-12 (16/36/64); analysis/calibration_vs_recorded.csv (calibrate_authors_code.py) -> 2 of 24 P < 0.05, 0.05 x 24 = 1.2. Lead: move into generated macros + check (R5-U3) if time permits
% CHECK-DIR [F31]: PyWake NOJ comparison at identical bins

\emph{Code, data and environment.} Runs: Linux containers (Intel Xeon; Python~3.11.15, NumPy~2.4.6, SciPy~1.17.1). Each run record gives the method, case, seed, budget, initialization, evaluations, final objective, feasibility, minimum spacing, coordinates and convergence curve; a guide names the command behind every run file, and one build script regenerates all macros, tables and figures and runs the checks. Code, per-run results and scripts are public (see the data availability statement).
% CHECK (manual, resolved 2026-09-29): versions pinned in analysis/requirements.txt ("from the container that produced the results"); clock speed and core count dropped because the analysis container (Xeon 2.1 GHz) and earlier run containers (2.8 GHz, final_prose.py) differ
% CHECK (manual, 2026-09-29): run-file columns Algorithm, Dataset, Radius, Turbines, Seed, Budget, Init, Calls, Objective, Ideal, WakeLoss, Feasible, MinSpacing, Seconds, Coordinates, Curve (mpce_*.csv); analysis/README_reproduce.md (command per run file); analysis/build_swevo.sh / build_mpce_paper.sh (mpce_results.py + check scripts)
